% 张三慧第4版第二卷《核物理》第30章思考题30.1--30.10
% V2主教材PDF物理页485（印刷468），不与同号章末习题混合计数。
\section*{第30章\quad 核物理：思考题}
\paragraph{思考题 30.1}
核子组成的原子核在何种近似下可想象为许多小球聚合在一起？
为什么不能把核当成经典硬小球的简单堆积？
\paragraph{思考题 30.2}
为什么不同原子核的平均核物质密度在粗略近似下几乎一样？
推导核半径随质量数 \(A\) 的尺度规律。
\paragraph{思考题 30.3}
核结合能主要由核子间的相互作用决定。
为什么总结合能随核子数 \(A\) 的主要项近似与 \(A\) 成正比？
说明核力饱和的作用。
\paragraph{思考题 30.4}
质子间库仑排斥很强，为什么仍能形成稳定原子核？
为什么过重原子核的稳定性降低？
\paragraph{思考题 30.5}
假设核电荷均匀分布在半径
\(R=r_0A^{1/3}\) 的球体内。
从均匀带电球的静电自能推出
半经验结合能公式中库仑修正系数 \(a_c\) 的近似表达式。
\paragraph{思考题 30.6}
完成原教材给出的放射变化并检查核素守恒：
\[
{}^{238}\mathrm U\to{}^{234}\mathrm{Th}+?,
\quad
{}^{90}\mathrm{Sr}\to{}^{90}\mathrm Y+?,
\]
\[
{}^{29}\mathrm{Cu}\to{}^{29}\mathrm{Ni}+?,
\quad
{}^{29}\mathrm{Cu}+?\to{}^{29}\mathrm{Zn}.
\]
注意原书最后两项对Cu、Ni、Zn的核电荷数/质量数标注存在排印歧义，
应先作核电荷守恒核查，而非猜填。
\paragraph{思考题 30.7}
从\({}^{235}\mathrm U\)放射系的起始核
至末态\({}^{207}\mathrm{Pb}\)，
一共经过多少次\(\alpha\)与\(\beta^-\)衰变？
由质量数与核电荷数守恒推导。
\paragraph{思考题 30.8}
为何通常自由原子核不能发生严格零反冲的
\(\gamma\)光子共振发射与吸收？
穆斯堡尔效应如何借助晶格实现近似无反冲的共振？
\paragraph{思考题 30.9}
为什么某些粒子—原子核反应截面
可以大于核的几何横截面？
举出量子波与共振吸收对截面的影响。
\paragraph{思考题 30.10}
为什么有些核反应虽放能
\(Q>0\)，仍需粒子具有高于 \(|Q|\) 或高于零的
入射能才能明显发生？
提高引发反应的粒子能量是否必然提高可利用净能量效率？
